\section{Experiments}\label{sec:experiments}

All experiments are reproducible from \lean{research/grh-pseudorandomness/experiments/} in the accompanying repository. For a prime $q$, the group $\Gq$
is cyclic of order $q-1$. After fixing a primitive root and the discrete-logarithm table, the biases against
all $q-1$ characters are the entries of a single length-$(q-1)$ FFT. The biases are therefore computed
exactly, up to floating-point rounding, not estimated by sampling. These numbers describe typical moduli
and are not proofs.

\paragraph{Primes versus random sets.}
Table~\ref{tab:typical} compares the uniform measure on primes $p\le x$ with a uniformly random multiset of
residues of the same size (\lean{prime_bias.py}). On typical prime moduli up to $10^6$, primes are as good as
random sets, and in fact slightly better. The normalized bias $\lambda_*\sqrt{\abs{\mathrm{supp}}/\log q}$ is
$0.85$--$0.95$, the Alon--Roichman scale. The GRH-provable scale $\log q/\sqrt x$ is pessimistic by roughly
$\sqrt{\log q}$ on typical moduli; adversarial moduli at accessible sizes did not change this
(\lean{adversarial_q.py}).

\begin{table}[ht]
\centering
\begin{tabular}{rrrr}
\toprule
$q$ & $x$ & $\lambda_*$, primes $\le x$ & $\lambda_*$, random multiset\\
\midrule
10007 & 3200 & 0.1182 & 0.1474\\
100003 & 3200 & 0.1378 & 0.1619\\
1000003 & 3200 & 0.1487 & 0.1631\\
\bottomrule
\end{tabular}
\caption{Exact bias of primes $p\le x$ against all nontrivial characters, versus a random multiset of the same
size ($452$).}\label{tab:typical}
\end{table}

\paragraph{The bounds dominate the truth.}
For the theorem's own measure we compared the exact bias with the theorem's bound. In
\lean{cosh_kernel_check.py} the bound uses the explicit-formula constants, evaluated by quadrature: the
measured bias is $0.50$--$0.74$ of the bound in $36/36$ cases ($q\le3.8\cdot10^6$, $a\in\{1,1.5,2\}$). In
\lean{sharp_check.py} the bound is $\mathrm{Num}/\mathrm{Den}$ of Theorem~\ref{thm:sharp} with quadrature
constants: the ratio is $0.03$--$0.27$ in $8/8$ cases. A violation in either case would have revealed an
error in the numerics, since the theorem itself is machine-checked. For the certified measure of
Corollary~\ref{cor:cert}, the exact bias for $q\in\{10007,100003,1000003,4000037\}$ is
$0.021$, $0.026$, $0.032$, $0.032$ (\lean{certificate_check.py}), against the worst-case guarantee $\frac12$. For the
smaller measure of Corollary~\ref{cor:convcert} it is $0.037$, $0.047$, $0.048$, $0.055$ (\lean{convsq_certificate_check.py}).

\paragraph{What the theorems certify.}
Table~\ref{tab:params} gives the support needed to certify bias $\le\varepsilon$ for all $q\le2^{64}$.

\begin{table}[ht]
\centering
\begin{tabular}{lrrr}
\toprule
 & $\varepsilon=0.5$ & $\varepsilon=0.1$ & $\varepsilon=0.01$\\
\midrule
Theorem~\ref{thm:A} via Corollary~\ref{cor:A} (formal) & $9.7\cdot10^{15}$ & $3.5\cdot10^{17}$ & $5.6\cdot10^{19}$\\
Theorem~\ref{thm:closed}, closed form (formal once instantiated) & $1.3\cdot10^{5}$ & $1.8\cdot10^{6}$ & $1.5\cdot10^{8}$\\
Corollary~\ref{cor:cert} (formal, instantiated) & $1.66\cdot10^{5}$ & --- & ---\\
Theorem~\ref{thm:convclosed}, convolution square (formal once instantiated) & $4.2\cdot10^{4}$ & $6.7\cdot10^{5}$ & $5.9\cdot10^{7}$\\
Corollary~\ref{cor:convcert} (formal, instantiated) & $4.84\cdot10^{4}$ & --- & ---\\
Theorem~\ref{thm:sharp}, constants by quadrature & $3.5\cdot10^{4}$ & $5.5\cdot10^{5}$ & $4.7\cdot10^{7}$\\
Leading term $c^{*2}\log^2q/\varepsilon^2$, positive-definite (Cor.~\ref{cor:asymp}) & $1.2\cdot10^{4}$ & $2.9\cdot10^{5}$ & $2.9\cdot10^{7}$\\
Leading term $\log^2q/\varepsilon^2$, barrier (Cor.~\ref{cor:asymp}) & $7.9\cdot10^{3}$ & $2.0\cdot10^{5}$ & $2.0\cdot10^{7}$\\
\bottomrule
\end{tabular}
\caption{Support $X$ (largest integer in the support) sufficient for bias $\le\varepsilon$ for every $q\le2^{64}$.
Rows two and four give the smallest support for which the closed-form ratio is $\le\varepsilon$, found numerically
(\lean{closed_form_params.py}, \lean{convsq_formal_params.py}); their numerical inequalities are machine-checked only for the
instances in rows three and five, which use round parameters ($y=2\log150$; $y=2\log220$, $R=12$). The last two rows omit
lower-order terms.}\label{tab:params}
\end{table}
